% 本文件由 LaTeX 排版 Agent 根据有效 translation.json 自动生成。
% author=Councils; work_id=3817; title=君士坦丁堡会议（主后394年）

\chapter{君士坦丁堡会议（主后394年）}

% structure: p0001 source=h1 target=omitted reason=page-title-already-rendered-by-parent

在我等最虔诚且深受天主爱戴的皇帝，弗拉维乌斯·阿卡狄乌斯·奥古斯都第三次任执政官，及霍诺留第二次任执政官期间，十月朔日前三日，在\Place{君士坦丁堡}至圣堂的洗礼堂中，当最神圣的主教们就座后\Emph{此处列有名单}，\Place{君士坦丁堡}主教\Person{内克塔留斯}说：\Quote{自从因天主的恩宠，本次会议在这神圣之地召开，如果我神圣的弟兄及圣事中的同工们认为妥当，因我看到我兄弟\Person{巴加迪乌斯}和\Person{阿加皮乌斯}也在场，他们正在为\Place{波斯特拉}的主教职位彼此争执，请让他们开始陈述各自的权利。}之后，他们为此事进行了一些程序，并表明前述\Person{巴加迪乌斯}仅由两位主教罢免，而这两位主教均已去世。\Place{安塞拉}主教\Person{阿拉比亚努斯}说：\Quote{并非因为这次判决，而是为我的余生担忧，我恳请神圣的会议制定一条法令：一位主教是否能够仅由两位主教罢免，无论都主教是否在场，且不影响当前案件。因为我担心，有人会从这些行为中获得权力，敢于尝试此类事情。因此，我希望得到你们的答复。}\par

\Place{君士坦丁堡}主教\Person{内克塔留斯}说：\Quote{最虔诚的主教\Person{阿拉比亚努斯}说得非常值得称道。但既然在判决上无法回溯，让我们在不谴责过去的情况下，为将来设立规则。}\Place{安塞拉}主教\Person{阿拉比亚努斯}说：\Quote{在\Place{尼西亚}召开的真福教父们的会议谴责了所发生的事，因为它命令不可少于三人祝圣，且即使如此，若没有都主教也不行。但是关于将来，我满怀恐惧提出了这个问题。因此，我希望你们清楚且毫不迟疑或怀疑地说，根据尼西亚会议的谕令，一位主教不能被两个人合法祝圣或罢免。}\par

之后，经过进一步辩论，\Place{亚历山大里亚}主教\Person{德奥菲卢斯}说：\Quote{对于已经离开的人，不能宣判任何愤怒的判决，因为应受谴责的人不在场。但如果有人考虑将来那些将要被罢免的人，在我看来，不仅这些人应该聚集，而且尽可能其他的省区主教也应聚集，以便通过众人的判决，对在场受审且应受罢免的人作出更准确的定罪。}\Place{君士坦丁堡}主教\Person{内克塔留斯}说：\Quote{既然争议涉及合法的制度和谕令，因此不可因个人原因而制定任何法令。因此，正如最神圣的主教\Person{阿拉比亚努斯}所言，为使将来确定，最神圣的主教\Person{德奥菲卢斯}的判决始终一致且审慎地规定：将来，甚至三位主教也不可罢免作为被告人接受审查的人，更何况两位主教，而应根据宗徒法典，由大会议和省区主教们判决。}\Place{安提约基雅}主教\Person{弗拉维安}说：\Quote{最神圣的主教\Person{内克塔留斯}和最神圣的主教\Person{德奥菲卢斯}所陈述的显然正确。}所有神职人员都同意这些。\par

\SourceRecord{Translated by Henry Percival.；Nicene and Post-Nicene Fathers, Second Series；Vol. 14.；Edited by Philip Schaff and Henry Wace.；Buffalo, NY: Christian Literature Publishing Co.,；1900.；<http://www.newadvent.org/fathers/3817.htm>.}
